\documentclass{article}

% if you need to pass options to natbib, use, e.g.:
%     \PassOptionsToPackage{numbers, compress}{natbib}
% before loading neurips_2026

% The authors should use one of these tracks.
% Before accepting by the NeurIPS conference, select one of the options below.
% 0. "default" for submission
\usepackage{neurips_2026}
% the "default" option is equal to the "main" option, which is used for the Main Track with double-blind reviewing.
% 1. "main" option is used for the Main Track
%  \usepackage[main]{neurips_2026}
% 2. "position" option is used for the Position Paper Track
%  \usepackage[position]{neurips_2026}
% 3. "eandd" option is used for the Evaluations & Datasets Track
 % \usepackage[eandd]{neurips_2026}
 % if you want to opt-in for a de-anomymized submission (not recommended, but OK):
 %\usepackage[eandd, nonanonymous]{neurips_2026}
% 4. "creativeai" option is used for the Creative AI Track
%  \usepackage[creativeai]{neurips_2026}
% 5. "sglblindworkshop" option is used for the Workshop with single-blind reviewing
 % \usepackage[sglblindworkshop]{neurips_2026}
% 6. "dblblindworkshop" option is used for the Workshop with double-blind reviewing
%  \usepackage[dblblindworkshop]{neurips_2026}

% After being accepted, the authors should add "final" behind the track to compile a camera-ready version.
% 1. Main Track
 % \usepackage[main, final]{neurips_2026}
% 2. Position Paper Track
%  \usepackage[position, final]{neurips_2026}
% 3. Evaluations & Datasets Track
 % \usepackage[eandd, final]{neurips_2026}
% 4. Creative AI Track
%  \usepackage[creativeai, final]{neurips_2026}
% 5. Workshop with single-blind reviewing
%  \usepackage[sglblindworkshop, final]{neurips_2026}
% 6. Workshop with double-blind reviewing
%  \usepackage[dblblindworkshop, final]{neurips_2026}
% Note. For the workshop paper template, both \title{} and \workshoptitle{} are required, with the former indicating the paper title shown in the title and the latter indicating the workshop title displayed in the footnote.
% For workshops (5., 6.), the authors should add the name of the workshop, "\workshoptitle" command is used to set the workshop title.
% \workshoptitle{WORKSHOP TITLE}

% "preprint" option is used for arXiv or other preprint submissions
 % \usepackage[preprint]{neurips_2026}

% to avoid loading the natbib package, add option nonatbib:
%    \usepackage[nonatbib]{neurips_2026}

\usepackage[utf8]{inputenc} % allow utf-8 input
\usepackage[T1]{fontenc}    % use 8-bit T1 fonts
\usepackage{hyperref}       % hyperlinks
\usepackage{url}            % simple URL typesetting
\usepackage{booktabs}       % professional-quality tables
\usepackage{amsfonts}       % blackboard math symbols
\usepackage{nicefrac}       % compact symbols for 1/2, etc.
\usepackage{microtype}      % microtypography
\usepackage{xcolor}         % colors

% Note. For the workshop paper template, both \title{} and \workshoptitle{} are required, with the former indicating the paper title shown in the title and the latter indicating the workshop title displayed in the footnote. 
\title{Formatting Instructions For NeurIPS 2026}


% The \author macro works with any number of authors. There are two commands
% used to separate the names and addresses of multiple authors: \And and \AND.
%
% Using \And between authors leaves it to LaTeX to determine where to break the
% lines. Using \AND forces a line break at that point. So, if LaTeX puts 3 of 4
% authors names on the first line, and the last on the second line, try using
% \AND instead of \And before the third author name.


\author{%
  David S.~Hippocampus\thanks{Use footnote for providing further information
    about author (webpage, alternative address)---\emph{not} for acknowledging
    funding agencies.} \\
  Department of Computer Science\\
  Cranberry-Lemon University\\
  Pittsburgh, PA 15213 \\
  \texttt{hippo@cs.cranberry-lemon.edu} \\
  % examples of more authors
  % \And
  % Coauthor \\
  % Affiliation \\
  % Address \\
  % \texttt{email} \\
  % \AND
  % Coauthor \\
  % Affiliation \\
  % Address \\
  % \texttt{email} \\
  % \And
  % Coauthor \\
  % Affiliation \\
  % Address \\
  % \texttt{email} \\
  % \And
  % Coauthor \\
  % Affiliation \\
  % Address \\
  % \texttt{email} \\
}


\begin{document}


\maketitle


\begin{abstract}

\end{abstract}



\section{Introduction}


\section{Related Work}
\subsection{Object Removal Datasets}
AvED-Bench
\subsection{Single-Modality Editing}
Prior work on audio and video editing largely treats the two modalities separately. For audio removal and editing, existing methods condition on text or audio cues without using aligned visual context: text-based diffusion approaches such as AUDIT~\cite{wang2023auditaudioeditingfollowing} and ZEUS~\cite{manor2024zeroshotunsupervisedtextbasedaudio} follow textual instructions, while audio-conditioned methods rely on reference audio or learned latent factors~\cite{cheng2025audio,liang2025audiomorphix}. These systems can remove or replace target sounds but ignore visual cues that help disambiguate on- and off-screen sources and enforce temporal alignment. On the visual side, video editors either adapt text-to-image or text-to-video diffusion backbones for object removal, or train on paired datasets. LGVI~\cite{wu2024languagedrivenvideoinpaintingmultimodal}, VideoPainter~\cite{bian2025videopainteranylengthvideoinpainting}, and VACE~\cite{jiang2025vaceallinonevideocreation} adapt text-to-video or video-diffusion backbones for object-level editing by conditioning on prompts and structured controls such as masks and text, while paired-supervision methods such as ROVI~\cite{wu2024languagedrivenvideoinpaintingmultimodal} learn directly from before/after examples. As shown in InstructPix2Pix~\cite{brooks2023instructpix2pixlearningfollowimage}, natural language instructions can serve as an expressive editing interface, but most video editors with text-to-video backbone rely on dense output captions. We adopt the paired-data paradigm and extend it to synchronized, object-level audiovisual edits with user-specified instructions.


ROSE\citep{miao2025roseremoveobjectseffects} uses mask as prompt

Effecterase\citep{fu2026effecterasejointvideoobject} uses mask as prompt

\subsection{Audiovisual Joint Generation}
Earlier audio–visual work primarily focuses on generating a single modality conditioned on the other \cite{sung2023sound, li2022learning, iashin2021taming, pascual2024masked, su2023physics, wang2024v2a, yang2024draw, liang2024foundations}. More recent efforts shift toward joint generation or editing of an audio–video pair, typically following two paradigms. In cascaded approaches, one generates video first and then synthesizes audio conditioned on it \cite{comunita2024syncfusion, hu2024video, jeong2025read}, while the other generates audio first and then synthesizes video conditioned on the sound \cite{jeong2023power, yariv2024diverse, zhang2024audio}. This design is simple and modular but can accumulate errors and drift in cross-modal coherence. In non-cascaded approaches, audio and video are generated jointly with objectives for synchronization, semantic consistency, and temporal alignment, yielding tighter coupling; recent work shows strong correspondence and spatiotemporal alignment \cite{zhao2025uniform, liu2024syncflow, tang2023any}.
However, most joint editing methods rely on diffusion inversion~\cite{liang2024language,lin2025zero} and text-to-video backbones, which typically restrict outputs to low frame rates~\cite{henschel2025streamingt2vconsistentdynamicextendable}. In this work, we synthesize audiovisual pairs and train on generated pairs directly to perform joint audiovisual removal.

\section{SAVEBench Dataset}
We use UltraVideo \citep{xue2025ultravideohighqualityuhdvideo}, VGGSound\citep{chen2020vggsoundlargescaleaudiovisualdataset} and Audioset\citep{gemmeke2017audioset} to construct 100K samples with the following pipeline. We apply several filtering method to filter out those bad samples

\subsection{Data Construction Pipeline}

\paragraph{Filtering}
We remove those samples with speech and background music with FireRedASR2S\citep{xu2026fireredasr2sstateoftheartindustrialgradeallinone}
for better semantic level correspondence


\paragraph{Target Object Identification}
We generate audiovisual captions for each sounding video with Qwen3-Omni\citep{xu2025qwen3omnitechnicalreport}. Then we prompt GPT-4o-mini to extract the object from the generated captions\citep{openai2024gpt4technicalreport}. We run sam3\citep{carion2025sam3segmentconcepts} on the first frame of each video and only keep those samples with more than 15\% object coverage

\paragraph{Target Object Track Separation}

We use sam-audio\citep{shi2025samaudiosegmentaudio}. We first use video mask to get target and residual. Then we use the object name to get the target and residual from the first round residual. 

\paragraph{Video Pair Construction}
We leverage Effecterase\citep{fu2026effecterasejointvideoobject} to remove the object on visual part. After we run the inpainting process, we run sam3 on the original video and removed video again and compare the result, we only keep those with 80\% mask left. 

\paragraph{Audiovisual Pairs}


\subsection{Data Statistics}
With our constructed dataset generation pipeline, we present object-level annotations with text and mask for our dataset

\section{SBFM}

\subsection{Audio-visual Removal}


\subsection{Flow Matching}


\subsection{Conditions}
We run with PE-Encoder\citep{bolya2025perceptionencoderbestvisual} where it is trained with audiovisual contrastive pairs. 

\section{Experiment Setup}
\subsection{Implementation Details}
\subsection{Baselines}
We compare our model against two baseline types: one is cross-modal baselines, where we cross two audio editors with three video editors, and the other is a multimodal editor trained with a conditional flow matching objective. 

Audiovisual Joint Editing\citep{lin2025zeroshotaudiovisualeditingcrossmodal}: AvED


\subsection{Evaluation Metrics}

\section{Results}
\subsection{Qualitative Examples}
\subsection{Quantitiave Results}
\section{Ablations}
\subsection{The bridge}

{
\small
\bibliographystyle{plainnat}
\bibliography{neurips_2026}
}

%%%%%%%%%%%%%%%%%%%%%%%%%%%%%%%%%%%%%%%%%%%%%%%%%%%%%%%%%%%%

\appendix

\section{Technical appendices and supplementary material}
Technical appendices with additional results, figures, graphs, and proofs may be submitted with the paper submission before the full submission deadline (see above). You can upload a ZIP file for videos or code, but do not upload a separate PDF file for the appendix. There is no page limit for the technical appendices. 

Note: Think of the appendix as ``optional reading'' for reviewers. The paper must be able to stand alone without the appendix; for example, adding critical experiments that support the main claims to an appendix is inappropriate. 

%%%%%%%%%%%%%%%%%%%%%%%%%%%%%%%%%%%%%%%%%%%%%%%%%%%%%%%%%%%%

\newpage
\input{checklist.tex}


\end{document}